% GENERATED FILE. Edit canonical lessons, then run npm run generate:books.
% canonical-source-hash: fnv1a64:7cd59a62ba09643d
% canonical-lessons: GE-C394-hochladen, GE-C394-webseite, GE-C394-suchmaschine, GE-C394-posten, GE-C394-kommentieren

\chapter{Im Internet}
\label{ch:ge-im-internet}
\hlchaptermodality{394}

\begin{chapteropening}
\textbf{By the end of this chapter:} \emph{I can talk about uploading, websites, search engines, posting and commenting.}

Complete the last lesson of chapter 394: I can talk about uploading, websites, search engines, posting and commenting.
\end{chapteropening}

\section[hochladen]{hochladen — to upload}

\label{lesson:GE-C394-hochladen}



\begin{quote}
Before the new one: say the German for the attachment, then the German for to attach.
\end{quote}

\subsection*{You'll want to know: hochladen}
\textbf{hochladen} — \textquotedblleft{}to upload\textquotedblright{}.

\begin{quote}
Sie hat das Video gestern hochgeladen.
\end{quote}

\begin{quote}
\emph{She uploaded the video yesterday.}
\end{quote}

Separable and strong: \emph{sie lädt hoch}, \emph{sie hat hochgeladen}.

\begin{cousinweb}
\emph{hochladen} renders English \emph{upload} part by part: \emph{hoch} \textquotedblleft{}up\textquotedblright{}, Old High German \emph{hōh} (EtymDB), and \emph{laden} \textquotedblleft{}to load\textquotedblright{}, Old High German \emph{hladan} (EtymDB).
\end{cousinweb}

\begin{grammarlens}[title={What you've built}]
One more word for telling, explaining and describing at B1.
\end{grammarlens}

\subsection*{Your turn}
\begin{itemize}
\raggedright
  \item \emph{Say it:} \emph{hochladen}
  \item \emph{Say it:} \emph{hochladen}, once more
  \item \emph{Say it:} a sentence of your own with \emph{hochladen}
  \item \emph{Recall:} say the German for the attachment, then the German for to attach, then say \emph{hochladen} again
\end{itemize}

\begin{tcolorbox}[breakable,colback=teal!4,colframe=teal!35!black,title={Before you move on}]
What does hochladen mean? (\textquotedblleft{}To upload\textquotedblright{}.) Say it once more.
\end{tcolorbox}
\section[die Webseite]{die Webseite — the web page, the website}

\label{lesson:GE-C394-webseite}



\begin{quote}
Before the new one: say the German for to attach, then the German for to upload.
\end{quote}

\subsection*{You'll want to know: die Webseite}
\textbf{die Webseite} — \textquotedblleft{}the web page, the website\textquotedblright{}.

\begin{quote}
Die Öffnungszeiten stehen auf der Webseite.
\end{quote}

\begin{quote}
\emph{The opening hours are on the website.}
\end{quote}

Plural: \emph{die Webseiten}. Also written \emph{die Website}, said the English way.

\begin{cousinweb}
\emph{Webseite} joins English \emph{web} to \emph{Seite} \textquotedblleft{}page, side\textquotedblright{}, inherited from Old High German \emph{sīta} and Proto-Germanic \emph{sīdōn} (EtymDB), after English \emph{web page}.
\end{cousinweb}

\begin{grammarlens}[title={What you've built}]
One more word for telling, explaining and describing at B1.
\end{grammarlens}

\subsection*{Your turn}
\begin{itemize}
\raggedright
  \item \emph{Say it:} \emph{die Webseite}
  \item \emph{Say it:} \emph{die Webseite}, once more
  \item \emph{Say it:} a sentence of your own with \emph{die Webseite}
  \item \emph{Recall:} say the German for to attach, then the German for to upload, then say \emph{die Webseite} again
\end{itemize}

\begin{tcolorbox}[breakable,colback=teal!4,colframe=teal!35!black,title={Before you move on}]
What does die Webseite mean? (\textquotedblleft{}The web page, the website\textquotedblright{}.) Say it once more.
\end{tcolorbox}
\section[die Suchmaschine]{die Suchmaschine — the search engine}

\label{lesson:GE-C394-suchmaschine}



\begin{quote}
Before the new one: say the German for to upload, then the German for the web page.
\end{quote}

\subsection*{You'll want to know: die Suchmaschine}
\textbf{die Suchmaschine} — \textquotedblleft{}the search engine\textquotedblright{}.

\begin{quote}
Mit einer Suchmaschine habe ich das Rezept schnell gefunden.
\end{quote}

\begin{quote}
\emph{I quickly found the recipe with a search engine.}
\end{quote}

Plural: \emph{die Suchmaschinen}.

\begin{cousinweb}
\emph{Suchmaschine} joins \emph{suchen} \textquotedblleft{}to search\textquotedblright{}, Proto-Germanic \emph{sōkijaną} (EtymDB), to \emph{Maschine} \textquotedblleft{}machine\textquotedblright{}, from Latin \emph{māchina} and Greek \emph{mēkhan\'{\={e}}} (EtymDB; Lewis \& Short).
\end{cousinweb}

\begin{grammarlens}[title={What you've built}]
One more word for telling, explaining and describing at B1.
\end{grammarlens}

\subsection*{Your turn}
\begin{itemize}
\raggedright
  \item \emph{Say it:} \emph{die Suchmaschine}
  \item \emph{Say it:} \emph{die Suchmaschine}, once more
  \item \emph{Say it:} a sentence of your own with \emph{die Suchmaschine}
  \item \emph{Recall:} say the German for to upload, then the German for the web page, then say \emph{die Suchmaschine} again
\end{itemize}

\begin{tcolorbox}[breakable,colback=teal!4,colframe=teal!35!black,title={Before you move on}]
What does die Suchmaschine mean? (\textquotedblleft{}The search engine\textquotedblright{}.) Say it once more.
\end{tcolorbox}
\section[posten]{posten — to post (online)}

\label{lesson:GE-C394-posten}



\begin{quote}
Before the new one: say the German for the web page, then the German for the search engine.
\end{quote}

\subsection*{You'll want to know: posten}
\textbf{posten} — \textquotedblleft{}to post (online)\textquotedblright{}.

\begin{quote}
Er postet jeden Tag Fotos von seinem Essen.
\end{quote}

\begin{quote}
\emph{He posts photos of his food every day.}
\end{quote}

Perfekt: \emph{hat gepostet}.

\begin{cousinweb}
\emph{posten} was borrowed from English \emph{post} (EtymDB). The English verb is usually traced to putting a notice up on a post, a word that goes back to Latin \emph{postis} \textquotedblleft{}a doorpost\textquotedblright{} (Lewis \& Short): you put a message up.
\end{cousinweb}

\begin{grammarlens}[title={What you've built}]
One more word for telling, explaining and describing at B1.
\end{grammarlens}

\subsection*{Your turn}
\begin{itemize}
\raggedright
  \item \emph{Say it:} \emph{posten}
  \item \emph{Say it:} \emph{posten}, once more
  \item \emph{Say it:} a sentence of your own with \emph{posten}
  \item \emph{Recall:} say the German for the web page, then the German for the search engine, then say \emph{posten} again
\end{itemize}

\begin{tcolorbox}[breakable,colback=teal!4,colframe=teal!35!black,title={Before you move on}]
What does posten mean? (\textquotedblleft{}To post (online)\textquotedblright{}.) Say it once more.
\end{tcolorbox}
\section[kommentieren]{kommentieren — to comment on}

\label{lesson:GE-C394-kommentieren}



\begin{quote}
Before the new one: say the German for the search engine, then the German for to post.
\end{quote}

\subsection*{You'll want to know: kommentieren}
\textbf{kommentieren} — \textquotedblleft{}to comment on\textquotedblright{}.

\begin{quote}
Viele Leute haben den Artikel kommentiert.
\end{quote}

\begin{quote}
\emph{Many people commented on the article.}
\end{quote}

It takes a direct object: \emph{einen Beitrag kommentieren}. \emph{Der Kommentar} — the comment.

\begin{cousinweb}
\emph{kommentieren} is a learned borrowing from Latin \emph{commentārī} \textquotedblleft{}to think over, to make notes on\textquotedblright{}, of the family of \emph{comminīscī} \textquotedblleft{}to devise, to invent\textquotedblright{} (EtymDB; Lewis \& Short).
\end{cousinweb}

\begin{grammarlens}[title={What you've built}]
One more word for telling, explaining and describing at B1.
\end{grammarlens}

\subsection*{Your turn}
\begin{itemize}
\raggedright
  \item \emph{Say it:} \emph{kommentieren}
  \item \emph{Say it:} \emph{kommentieren}, once more
  \item \emph{Say it:} a sentence of your own with \emph{kommentieren}
  \item \emph{Recall:} say the German for the search engine, then the German for to post, then say \emph{kommentieren} again
\end{itemize}

\begin{tcolorbox}[breakable,colback=teal!4,colframe=teal!35!black,title={Before you move on}]
What does kommentieren mean? (\textquotedblleft{}To comment on\textquotedblright{}.) Say it once more.
\end{tcolorbox}
